\subsection{Riesz 空间由其正元素生成}

设 \(\mathbf{E}\) 是一个序向量空间；集 \(\mathrm{P}\) （由元素 \(\geqslant 0\) 组成，元素取自 \(\mathbf{E}\) ）是一个以 0 为顶点的凸锥，也就是说 (TVS, II, §2, No. 4)，一个满足

\(P + P \subset P\) 且 \(\lambda P \subset P\) （对一切 \(\lambda > 0\) ）的集。反过来，若在 R 上的向量空间 E 中，P 是一个以 0 为顶点的凸锥，满足 \(P \cap (-P) = \{0\}\) （换句话说，一个带顶点的真凸锥），则已知 (loc. cit.) 关系 \(y - x \in P\) 是一个序关系（记作 \(x \leqslant y\) ），且与 E 的向量空间结构相容。为使这一序结构在 E 上定义一个 Riesz 空间结构，必要且充分的是：

\(1^{\circ}\) P 生成 E，也就是说，每个 \(z \in \mathbf{E}\) 都形如 \(y - x\) ，其中 \(x\) 与 \(y\) 属于 P；

\(2^{\circ}\) P 满足以下两个条件之一：

\begin{enumerate}
\def\labelenumi{\alph{enumi})}
\item
  \(\mathbf{P}\) 的每对元素在 \(\mathbf{P}\) 中有一个上确界；
\item
  P 的每对元素在 P 中有一个下确界 (A, VI, §1, No. 9, Prop. 8)。
\end{enumerate}

